% --------------------------------------------------------------------------------
% AutoRL_Combined_Final.tex
% Final combined manuscript (Article-2 base + Article-1 technical enrichments)
% Date: 19-Mar-2026
% --------------------------------------------------------------------------------

\documentclass[11pt]{article}

\usepackage[margin=1in]{geometry}
\usepackage{amsmath, amssymb, amsfonts}
\usepackage{booktabs}
\usepackage{siunitx}
\usepackage[table]{xcolor}
\usepackage{graphicx}
\graphicspath{{figures/}{paper/figures/}}
\usepackage{subcaption}
\usepackage{caption}
\usepackage{array}
\usepackage{setspace}
\usepackage{csquotes}
\usepackage{listings}
\usepackage{hyperref}
\usepackage[nameinlink,noabbrev]{cleveref}

\usepackage[style=apa,backend=biber,sorting=nyt]{biblatex}
\DeclareLanguageMapping{american}{american-apa}
\IfFileExists{references.bib}{\addbibresource{references.bib}}{\addbibresource{../references.bib}}

\newcommand{\autorlFigRoot}{\detokenize{figures and reports}}
\IfFileExists{\detokenize{paper/figures and reports/SIT_Results/_Analysis_Report_SIT_Set-1.jpg}}{%
  \renewcommand{\autorlFigRoot}{\detokenize{paper/figures and reports}}%
}{}

\newcommand{\autorlTableRoot}{tables}
\IfFileExists{paper/tables/hypothesis_tests_summary.tex}{%
	\renewcommand{\autorlTableRoot}{paper/tables}%
}{}

\hypersetup{
  colorlinks=true,
  linkcolor=blue,
  citecolor=blue,
  urlcolor=blue
}

\captionsetup{labelfont=bf, textfont=bf}
\sisetup{round-mode=places, round-precision=3, detect-weight=true}
\onehalfspacing

\lstset{
	basicstyle=\ttfamily\footnotesize,
	columns=fullflexible,
	breaklines=true,
	frame=single,
	xleftmargin=0.5em,
	xrightmargin=0.5em,
	aboveskip=0.8em,
	belowskip=0.8em
}

\title{AutoRL for Predictive Maintenance on the Edge: \texorpdfstring{\newline}{ }A Comparative Study of REINFORCE with Attention Mechanisms for Milling Tool Replacement}

\author{Rajesh Siraskar\\
\small V.2.0\\
\small (SCRI SIT Pune)\\
\small \texttt{(email)}
}

\date{\today}

\begin{document}
\maketitle

\begin{abstract}
Industry~4.0 maintenance systems increasingly rely on edge-computing deployments where stringent limits on memory, latency, and energy consumption favor computationally light learning methods. In this work we study \emph{AutoRL}, an automated reinforcement learning (RL) pipeline that trains and evaluates maintenance agents for a milling machine tool-replacement task, with particular focus on whether the na\"{i}ve Monte-Carlo policy-gradient algorithm \textsc{REINFORCE} can be competitive with more advanced methods such as \textsc{PPO}. The contribution is threefold: (i) an end-to-end AutoRL training and multi-round evaluation framework, (ii) a Gymnasium-compatible predictive-maintenance environment that hides tool wear from observations to enforce inference from sensor signals alone, and (iii) a systematic comparison across four RL algorithms (\textsc{A2C}, \textsc{DQN}, \textsc{PPO}, \textsc{REINFORCE}) and four attention mechanisms (Nadaraya--Watson, Temporal, Multi-Head, Self-Attention) integrated as Stable-Baselines3 feature extractors.

Across 48 trained agents (3 training files $\times$ 4 algorithms $\times$ 4 attention mechanisms) and 20 evaluation rounds per configuration, tested on six unseen files, \textsc{REINFORCE} paired with temporal or multi-head attention is consistently among the top-performing configurations and is statistically superior to A2C and DQN baselines in unseen-only panels (one-tailed Welch $t$-tests, $\alpha=0.05$). Against PPO, the evidence is schema-dependent, and results support \textsc{REINFORCE} as a lightweight, competitive alternative when combined with attention-based inductive biases aligned with edge and sustainability constraints.
\end{abstract}

\textbf{Keywords:} Predictive maintenance; Industry~4.0; Edge computing; Reinforcement learning; \textsc{REINFORCE}; \textsc{PPO}; Attention mechanisms; Milling machine; Sustainable AI.

% -----------------------------------------------------------------------------
% Parent narrative (from Article_VSC.tex)
% -----------------------------------------------------------------------------
\section{Context: Industry~4.0/5.0, predictive maintenance, and edge constraints}
Predictive maintenance (PdM) is a foundational Industry~4.0 capability: instrumented assets continuously produce sensor telemetry and maintenance decisions are optimized to reduce unplanned downtime and extend component life. In practical deployments, PdM must increasingly operate at the \emph{edge}---near machines and shop-floor controllers---to meet latency, bandwidth, and reliability constraints, and to reduce operational risk associated with centralized compute.

Edge deployment imposes hard resource limits (compute cycles, memory footprint, and energy), motivating \emph{lightweight} ML/RL models and training procedures. These constraints directly relate to sustainability considerations: lower compute and shorter training runs reduce electricity consumption and associated carbon emissions, which is increasingly relevant in corporate carbon accounting and Industry~5.0 human-centric and sustainable manufacturing initiatives. Recent work in ``Green AI'' emphasizes reporting and minimizing computational cost as a first-class metric alongside accuracy~\cite{Schwartz2020GreenAI, Strubell2019Energy}.

Beyond operational cost, reducing computational demand supports organizational carbon-footprint targets and can strengthen evidence bases used in carbon-credit and ESG reporting workflows, making ``efficient learning'' practically consequential rather than purely academic.

In this setting, a key research question emerges:
\begin{quote}
\emph{Can simple, computationally light RL algorithms such as \textsc{REINFORCE} be competitive with more advanced policy-optimization methods (e.g., \textsc{PPO}) for PdM decision-making, especially when augmented with attention-based feature extraction?}
\end{quote}

This study builds directly on prior findings that na\"{i}ve \textsc{REINFORCE} can perform strongly for PdM in untuned regimes and that automated RL frameworks can help industrial practitioners explore algorithm--hyperparameter design spaces~\cite{Siraskar2025NaiveREINFORCE}. It also leverages prior work applying a Nadaraya--Watson estimator as a lightweight attention mechanism for PdM, motivated by edge constraints~\cite{Siraskar2024NWAttention}.

\section{AutoRL overview (high level)}
AutoRL is designed as a practical experimentation platform to (i) train multiple RL agents over a grid of algorithmic choices and hyperparameters, (ii) evaluate agents on multiple test files using prognostic-aware metrics, and (iii) generate standardized analysis artefacts (dashboards, heatmaps, and hypothesis-test summaries) for model selection. In this repository snapshot, the primary entry point is a CLI pipeline (\texttt{code/batch\_train.py}) that supports training, discovery of trained models, and evaluation in single-round or multi-round modes. Multi-round evaluation produces a ``long'' CSV with all repetitions to support statistical analysis (e.g., Welch $t$-tests).

% -----------------------------------------------------------------------------
% Merged manuscript body (from AutoRL_for_PdM.tex)
% Note: Results section is preserved from AutoRL_for_PdM.tex.
% -----------------------------------------------------------------------------
\section{Introduction}
\label{sec:intro}
Predictive maintenance (PdM) aims to intervene \emph{before} functional failure while maximizing asset utilization and protecting product quality. In Industry~4.0 deployments, sensorized manufacturing assets generate rich telemetry streams that enable condition-aware decision support and prescriptive control \parencite{Zonta2020PdM, Compare2020IoT}. However, the operational setting for PdM is increasingly edge-centric: decisions must be made close to equipment with bounded latency and without guaranteed cloud connectivity. Under these constraints, the cost profile of learning methods matters. In particular, compute-intensive deep RL algorithms may be difficult to train broadly (e.g., within AutoRL sweeps) or to deploy reliably on constrained devices \parencite{Njor2024TinyML, Pandey2023EdgeDNN}.

RL is attractive for PdM because maintenance is not only a prediction problem; it is a sequential decision problem under uncertainty with asymmetric costs for late actions (failures) versus early actions (premature replacement). This paper studies PdM of a milling machine tool as a binary decision problem (continue vs replace) and focuses on a core research tension: whether a computationally light policy-gradient method (REINFORCE) can be competitive with (or better than) more advanced baselines when the environment is partially observable and depends on effective feature extraction from multi-sensor streams.

\subsection{Research objectives and hypotheses}
The study addresses the following research objectives (RO): (RO1) compare RL algorithms suitable for PdM, (RO2) implement an AutoRL pipeline to automatically generate RL agents for PdM, (RO3) design for industrial practitioners by reducing the need to understand RL or hyperparameter tuning, (RO4--RO5) design a Gymnasium-compatible RL environment that generalizes across sensor schemas, and (RO6--RO7) study edge constraints (and related sustainability motivations) and attention mechanisms for improving lightweight RL.

We test three hypotheses:
\begin{itemize}
	\item \textbf{H1 (primary):} A computationally light REINFORCE agent outperforms A2C, DQN, and PPO on the studied PdM task.
	\item \textbf{H2:} The same environment and reward design generalize across two different milling-machine schemas (SIT and IEEE).
	\item \textbf{H3:} Attention mechanisms improve REINFORCE sufficiently to match or exceed PPO, with multi-head attention providing the most robust gains.
\end{itemize}

\subsection{Contributions}
\label{sec:contrib}
This paper contributes:
\begin{itemize}
	\item \textbf{AutoRL pipeline for practitioners} (\cref{sec:autorl}): batch training and evaluation with checkpointed recovery to support large experiment grids.
	\item \textbf{A schema-adaptive Gymnasium environment} (\cref{sec:env}): a milling-tool PdM environment that hides wear from observations, enabling realistic partial observability across different sensor schemas \parencite{Gymnasium2023}.
	\item \textbf{Attention mechanisms implemented as lightweight feature extractors} (\cref{sec:attention}): Nadaraya--Watson, temporal attention, multi-head attention, and self-attention integrated into RL policies.
	\item \textbf{Empirical evaluation across SIT and IEEE schemas} (\cref{sec:results}) with multi-round evaluation on unseen data and hypothesis testing at $\alpha=0.05$.
\end{itemize}

\section{Background and related work}
\label{sec:background}
\subsection{PdM, Industry~4.0, and edge constraints}
PdM is commonly positioned as an Industry~4.0 centerpiece because it links sensing, connectivity, analytics, and operational decision-making in cyber-physical production systems \parencite{Compare2020IoT, Zonta2020PdM}. Prescriptive maintenance architectures further motivate this role by embedding decision-making (when to defer vs intervene) into cyber-physical workflows \parencite{Ansari2019PriMa}. Edge deployment strengthens the PdM case by enabling low-latency, resilient operation near equipment, but it also introduces feasibility constraints: memory, compute, latency, and system-integration overhead can limit the practical adoption of complex ML/RL pipelines \parencite{Njor2024TinyML, Pandey2023EdgeDNN}. These constraints connect to sustainability and ``Green AI'' concerns: when many candidate agents are trained during AutoRL exploration, the incremental cost per training run becomes operationally and environmentally relevant \parencite{Schwartz2020GreenAI, Strubell2019Energy}.

\subsection{RL for predictive maintenance}
RL has been surveyed extensively in PdM, with evidence that a wide range of formulations have been explored across maintenance scheduling, inspection, and resource-constrained decision problems \parencite{Siraskar2023RLReview, Ogunfowora2023, Zhang2024}. A consistent theme across these surveys is that the ``right'' RL complexity level depends on (i) the structure of the decision problem (action-space size, horizon, coupling) and (ii) deployment realism (compute, memory, latency).

%% From UM_BACKGROUND:
The literature also suggests a useful two-stream picture. Many PdM and condition-based maintenance studies learn effective policies using non-deep or otherwise lightweight RL, especially when state representations are structured and action spaces are modest \parencite{Adsule2020CBMRL, Yousefi2020QualityEngineering, Rasay2024JointOptimization, Dadash2024Lifespan}. Deep RL (e.g., DQN/PPO-style methods) becomes more prominent as environments become richer and harder: higher-dimensional sensing, long-horizon tradeoffs, multi-component scheduling, or coupled resource constraints \parencite{Liu2020EJOR, Chen2023IISE, LeeMitici2022RESS, Ong2022TransferDRL}.

This motivates a direct comparison between a lightweight method (REINFORCE) and widely-used deep RL baselines (A2C, DQN, PPO) under a shared environment and evaluation protocol.

%% From UM_literature_resarch:
\subsection{Positioning lightweight RL for edge PdM (qualified defense)}
The strongest literature-backed position is a \emph{qualified defense} of lightweight RL for edge PdM: the literature does not establish that REINFORCE is universally better than advanced deep RL, but it does support the claim that lightweight methods are often the more credible engineering choice when the PdM task is structured and edge constraints are first-order requirements \parencite{Compare2020IoT, Njor2024TinyML, Pandey2023EdgeDNN}. In embedded RL more broadly, deep RL frequently requires additional machinery---transfer learning, pruning, compression, or distillation---to fit acutely constrained devices \parencite{Jang2020KnowledgeTransfer, Szydlo2022TinyRL, Xu2022ISORC}.

Sustainability motivations should also be framed carefully. Maintenance optimization can contribute to energy and emissions objectives in industrial systems, but the PdM literature rarely reports algorithm-level energy draw or carbon footprints for competing RL methods. Accordingly, ``Green AI'' claims in this paper are positioned as downstream implications of reduced computational demand and edge-feasible deployment, not as directly measured environmental outcomes \parencite{Banyai2021Energies, Banyai2022Sustainability, JasiulewiczKaczmarek2024Maintenance40}.

This framing also shapes how we interpret statistical comparisons. Hypothesis tests in this paper are used to argue that observed differences are unlikely under the chosen experimental design, not to claim universal superiority across all PdM deployments \parencite{Zhang2024}.

\subsection{AutoRL and practitioner-centered experimentation}
AutoRL refers to automating choices that otherwise require expert RL knowledge: algorithm selection, hyperparameter tuning, reward shaping, and robust evaluation. Prior work emphasizes that automation can improve performance and reproducibility, and can lower the barrier for non-expert adoption \parencite{Faust2019EvolvingRewards, Mussi2023ARLO}. For industrial practitioners, this design goal is particularly important: the tool should help them select a high-performing agent without requiring deep understanding of RL internals.

\subsection{Attention mechanisms for PdM}
Attention mechanisms reweight inputs to focus computation on the most relevant parts of an observation, and have become a standard architectural primitive in modern deep learning \parencite{Vaswani2017Attention}. In PdM, attention can be interpreted as a structured inductive bias that emphasizes salient sensor channels or degradation stages. The present study builds on prior PdM attention work that introduced a Nadaraya--Watson estimator based attention mechanism as a lightweight option aligned with edge constraints \parencite{Siraskar2024NWAttention}.

\subsection{Lightweight vs deep RL in PdM}
The PdM RL literature includes both lightweight and deep formulations. Many maintenance decision problems have small action spaces and can be solved with structured state representations and classical or lightweight RL (e.g., condition-based maintenance policies learned without large deep networks) \parencite{Adsule2020CBMRL, Yousefi2020QualityEngineering}. In contrast, deep RL is often justified when the environment is richer: larger state spaces, coupled resources, long horizons, or multi-component scheduling \parencite{Liu2020EJOR, LeeMitici2022RESS, Chen2023IISE, Ong2022TransferDRL}. This distinction is central to the present paper: the milling tool replacement task is intentionally formulated as a compact binary decision under partial observability, which makes it a good candidate for lightweight RL with carefully designed representations.

\subsection{Edge deployment, TinyML, and embedded RL}
Edge constraints are not merely an implementation detail; they influence which methods are feasible to train broadly (as in AutoRL sweeps) and to deploy reliably \parencite{Compare2020IoT, Njor2024TinyML, Pandey2023EdgeDNN}. Embedded RL research shows that deep RL is often paired with knowledge transfer, pruning, or distillation before it fits constrained devices \parencite{Jang2020KnowledgeTransfer, Szydlo2022TinyRL, Xu2022ISORC}. This motivates our emphasis on REINFORCE and lightweight attention feature extractors: if performance is statistically competitive, simplicity can become a defensible engineering preference rather than a compromise.

\section{Foundations: PdM as an MDP/POMDP and RL algorithms}
\label{sec:foundations}
\subsection{PdM as a decision process (MDP)}
\label{sec:mdp}
We model the milling-tool replacement task as an episodic Markov decision process (MDP)
$\mathcal{M}=\langle \mathcal{S},\mathcal{A},P,r,\gamma\rangle$ \parencite{SuttonBarto2018}.

\begin{itemize}
\item $\mathcal{S}$ is the (latent) state space capturing tool condition, operating parameters, and environment. We adopt the standard Markov assumption that the current state summarizes the history relevant for predicting future degradation.
\item $\mathcal{A}=\{0,1\}$ is a binary action space with $a=0$ meaning \emph{replace} and $a=1$ meaning \emph{continue}. This formulation is closely related to an optimal stopping decision: choose the replacement time that balances utilization against violation risk.
\item $P(s'\mid s,a)$ is the transition kernel.
\item $r(s,a)$ is a shaped reward encoding asymmetric PdM costs (\cref{sec:reward}).
\item $\gamma$ is the discount factor.
\end{itemize}

The agent observes state $s_t$ (or an observation $o_t$ in the POMDP case), samples an action $a_t\sim\pi_\theta(\cdot\mid s_t)$, and receives reward $r_t=r(s_t,a_t)$.

\subsection{Finite vs. infinite horizon and the discount factor $\gamma$}
\label{sec:gamma}
The admissible range of $\gamma$ depends on whether the process is finite or infinite horizon.

\paragraph{Infinite horizon ($\gamma\in[0,1)$).}
In non-terminating environments, $\gamma$ must be strictly less than $1$ to ensure the discounted return converges and the value function is well-defined.

\paragraph{Finite horizon ($\gamma\in[0,1]$).}
Predictive maintenance is naturally episodic: a rollout terminates when the tool is replaced or when a terminal failure/threshold event occurs. If the episode length is bounded by $T<\infty$, the cumulative return $\sum_{t=0}^{T}\gamma^t r_t$ is finite even when $\gamma=1$. In our AutoRL grid we therefore include $\gamma\in\{0.95,0.99,1.0\}$ as a scientifically valid choice in the finite-horizon formulation.

\subsection{Objective function, returns, and baselines}
\label{sec:objective}
The RL objective is to maximize expected discounted return under the policy $\pi_\theta$:
\begin{equation}
J(\theta)=\mathbb{E}_{\tau\sim\pi_\theta}\left[\sum_{t=0}^{T}\gamma^t r_t\right],
\end{equation}
where $\tau=(s_0,a_0,r_0,\ldots,s_T)$ is a trajectory.
Define the return from time $t$ as
\begin{equation}
G_t = \sum_{k=t}^{T} \gamma^{k-t} r_k.
\end{equation}
Policy-gradient methods often subtract a baseline $b(s_t)$ (e.g., a value function) from $G_t$ to reduce variance without changing the expected gradient.

\subsection{Partial observability in milling PdM (POMDP)}
\label{sec:pomdp}
In realistic PdM, the true health state (tool wear) is not directly observable at runtime. The environment therefore corresponds to a partially observable MDP (POMDP), where the agent receives an observation vector $o_t$ derived only from sensor channels and must infer the hidden wear variable $w_t$:
\begin{equation}
	o_t = O(s_t) = \mathbf{x}_t, \qquad w_t \notin o_t.
\end{equation}
This design choice is implemented explicitly in the environment (\cref{sec:env}), and motivates attention-based feature extraction (\cref{sec:attention}) to improve inference from sensor streams.

\subsection{RL algorithms compared}
\label{sec:algos}
We compare four algorithms: DQN (value-based), A2C (actor--critic), PPO (policy optimization), and REINFORCE (Monte-Carlo policy gradient). We summarize their update principles and provide the key objectives used in practice.

\paragraph{REINFORCE (Monte-Carlo policy gradient).}
REINFORCE directly estimates the policy gradient using Monte-Carlo returns \parencite{Williams1992REINFORCE}:
\begin{equation}
\nabla_\theta J(\theta) = \mathbb{E}\left[\sum_{t=0}^{T} \nabla_\theta \log \pi_\theta(a_t\mid s_t)\,\big(G_t-b(s_t)\big)\right].
\end{equation}
Its simplicity makes it attractive in AutoRL sweeps; the primary challenge is variance, which motivates baselines and/or representation learning.

\paragraph{A2C (advantage actor--critic).}
Actor--critic methods use a learned value baseline $V_\phi(s_t)$ to reduce variance. With advantage estimate $\hat{A}_t=G_t-V_\phi(s_t)$, the policy gradient takes the form
\begin{equation}
\nabla_\theta J(\theta) \approx \mathbb{E}\big[\nabla_\theta\log\pi_\theta(a_t\mid s_t)\,\hat{A}_t\big].
\end{equation}
The critic is trained by minimizing a regression loss such as
\begin{equation}
\mathcal{L}_V(\phi)=\mathbb{E}\big[(G_t - V_\phi(s_t))^2\big].
\end{equation}

\paragraph{PPO (clipped policy optimization).}
PPO constrains policy updates using a clipped surrogate objective \parencite{Schulman2017PPO}. Define the probability ratio
\begin{equation}
 r_t(\theta)=\frac{\pi_\theta(a_t\mid s_t)}{\pi_{\theta_{\mathrm{old}}}(a_t\mid s_t)}.
\end{equation}
With advantage $\hat{A}_t$, the clipped objective is
\begin{equation}
\mathcal{L}^{\mathrm{CLIP}}(\theta)=\mathbb{E}\left[\min\Big(r_t(\theta)\hat{A}_t,\;\mathrm{clip}(r_t(\theta),1-\epsilon,1+\epsilon)\hat{A}_t\Big)\right].
\end{equation}
In typical implementations, the optimized objective also includes a value-function penalty and an entropy bonus:
\begin{equation}
\mathcal{J}(\theta,\phi)=\mathcal{L}^{\mathrm{CLIP}}(\theta) - c_1\,\mathcal{L}_V(\phi) + c_2\,H(\pi_\theta).
\end{equation}

\paragraph{DQN (deep Q-learning).}
DQN learns an action-value function $Q_\theta(s,a)$ and minimizes a temporal-difference error with a target network \parencite{Mnih2015DQN}:
\begin{equation}
\mathcal{L}(\theta)=\mathbb{E}\left[\left(r_t + \gamma\max_{a'}Q_{\bar\theta}(s_{t+1},a') - Q_\theta(s_t,a_t)\right)^2\right].
\end{equation}
Experience replay and target-network updates stabilize learning, but partial observability can degrade performance without explicit memory.

\paragraph{Why REINFORCE can be edge-suited.}
In an AutoRL setting, algorithmic simplicity can translate to lower total compute when sweeping many configurations. This does not guarantee superior reward, but it shifts the burden of proof: if a lightweight method is statistically competitive on PdM metrics, it may be preferable for edge-first deployments \parencite{Schwartz2020GreenAI, Njor2024TinyML}.

\section{Problem formulation: asymmetric costs and decision rationale}
\label{sec:problem}
PdM decisions are inherently asymmetric. A late replacement can cause defective parts, scrap, or catastrophic failure, while an early replacement wastes remaining tool life. This asymmetry motivates a policy that is conservative near the wear threshold, yet aims to maximize utilization.

\subsection{A cost-aware view of the binary action}
Let $\tau$ denote the true wear threshold (normalized to $w=1.0$). At each timestep $t$ the agent chooses continue or replace. Although the implemented reward is not a direct monetary cost model, it is designed to approximate these asymmetries by combining (i) a survival incentive for safe continuation, (ii) a catastrophic penalty for threshold violation, and (iii) a bonus for replacing close to the threshold.

\begin{table}[t]
	\centering
	\caption{Conceptual cost-aware interpretation of PdM decisions (qualitative), aligned with the reward shaping in \cref{sec:reward}.}
	\label{tab:cost_view}
	\begin{tabular}{p{0.18\linewidth} p{0.32\linewidth} p{0.42\linewidth}}
		\toprule
		\rowcolor{gray!20} \textbf{Case} & \textbf{Condition} & \textbf{PdM interpretation}\\
		\midrule
		Timely replace & $a_t=0$ near $w\approx\tau$ & Desired: high utilization and low violation risk.\\
		Early replace & $a_t=0$ when $w\ll\tau$ & Safe but wasteful: reduces risk while discarding useful life.\\
		Late action / violation & $a_t=1$ while $w>\tau$ & Worst case: elevated risk of quality loss or failure; strongly penalized.\\
		Continue safely & $a_t=1$ with $w<\tau$ & Encouraged: extracts remaining useful life when safe.\\
		\bottomrule
	\end{tabular}
\end{table}

\subsection{Why attention matters under partial observability}
Wear is hidden from the observation (\cref{sec:pomdp}). The agent must therefore infer degradation from noisy sensor vectors and their evolution. In such a POMDP, representation quality can dominate algorithmic sophistication: a lightweight policy-gradient method can succeed if the feature extractor emphasizes wear-correlated channels and suppresses nuisance variation.

\section{AutoRL pipeline for practitioner-centered experimentation}
\label{sec:autorl}
The AutoRL pipeline is implemented as a CLI batch runner (\texttt{code/batch\_train.py}) that supports training, evaluation, and reporting across a grid of algorithms, hyperparameters, and attention mechanisms. The pipeline is designed to be practitioner-friendly in two ways:
\begin{itemize}
	\item \textbf{Automation:} algorithm and attention choices are enumerated automatically; evaluation is repeated for multiple rounds.
	\item \textbf{Reliability:} long-running sweeps use a lightweight SQLite checkpoint queue to recover from failures without manual intervention.
\end{itemize}
For this study, multi-round evaluation uses 20 rounds (\texttt{EVAL\_ROUNDS=20}) to support confidence intervals and statistical tests.

\subsection{Workflow, outputs, and reproducibility artefacts}
The CLI provides: (i) training and evaluation orchestration, (ii) model discovery, and (iii) plot/report generation. Typical artefacts include trained model files, per-run metadata JSON (algorithm, attention type, training file, learning rate, $\gamma$, episodes), ``long'' CSV outputs for multi-round evaluation, and summary plots (analysis dashboards, generalization heatmaps, and statistical heatmaps). Checkpointing is implemented via a lightweight SQLite database to support recovery from interruptions.

\subsection{Search space and experiment grid}
The AutoRL study evaluates combinations of:
\begin{itemize}
	\item \textbf{Algorithm:} A2C, DQN, PPO, REINFORCE.
	\item \textbf{Attention:} None (baseline), Nadaraya--Watson (NW), temporal attention (TP), multi-head attention (MH), self-attention (SA).
	\item \textbf{Dataset schema:} SIT Set-1, SIT Set-2, IEEE.
\end{itemize}
Each configuration produces per-round evaluation metrics, enabling aggregation by algorithm and attention. This structure is important for practitioner-facing reporting: it supports robust selection (e.g., by unseen-only mean evaluation score) rather than anecdotal, single-run results.

\subsection{Checkpointing and failure recovery}
AutoRL sweeps can be long-running and failure-prone (e.g., interrupted runs, I/O issues). The pipeline therefore uses a lightweight SQLite-backed queue to record completed jobs and to retry failed configurations. This improves reproducibility and reduces manual intervention, which is critical when the end-user is an industrial practitioner rather than an RL specialist.

\section{Environment design: a Gymnasium-compatible milling PdM POMDP}
\label{sec:env}
\subsection{Design goals and schema abstraction}
The environment class \texttt{MT\_Env} (\texttt{code/rl\_pdm.py}) implements the Gymnasium interface \parencite{Gymnasium2023}. A central design goal is schema generality: the same environment should work across different milling-machine sensor schemas. This is implemented by (i) detecting the schema from column names and (ii) selecting the relevant sensor feature set accordingly.

\paragraph{Schema detection.}
The environment detects the schema using available sensor columns (e.g., force and acoustic emission channels for IEEE; vibration, sound, and load-cell channels for SIT). This supports H2 by enabling the same logic to run across the two datasets without changing the reward or transition logic.

\subsection{Observation and action spaces}
Only sensor readings are included in the observation vector; \emph{tool wear} is excluded (\cref{sec:pomdp}). The action space is discrete with two actions: replace or continue.

\subsection{Reward shaping for asymmetric PdM risk}
\label{sec:reward}
The reward function encourages long tool use without threshold violation, while rewarding timely replacement close to the wear threshold. In the implementation, continuing yields a small positive survival reward $R_1$, violating the wear threshold yields a catastrophic penalty $R_2$, replacement has a small cost $R_3$, and replacement near the threshold earns a strong bonus $R_4$ (with stepped bonuses above 70\%, 80\%, and 90\% wear ratio).

This design reflects a practical PdM asymmetry: a late replacement causing quality loss or failure is substantially worse than an early replacement.

\section{Attention mechanisms as lightweight feature extractors}
\label{sec:attention}
Because wear is hidden from observation, the core challenge is representation learning from sensor vectors. We implement attention mechanisms as feature extractors $\phi_\theta(\mathbf{x}_t)$ that map a raw sensor vector into a compact embedding used by the RL policy.

\paragraph{General attention concept.}
Given values $V=\{\mathbf{v}_1,\ldots,\mathbf{v}_n\}$ and a query $\mathbf{q}$, attention forms a convex combination
\begin{equation}
\mathrm{Attention}(\mathbf{q},V)=\sum_{i=1}^{n}\alpha_i\,\mathbf{v}_i,\qquad \sum_i\alpha_i=1,\;\alpha_i\ge 0,
\end{equation}
with weights computed by a softmax over a compatibility score $s(\mathbf{q},\mathbf{v}_i)$:
\begin{equation}
\alpha_i = \frac{\exp(s(\mathbf{q},\mathbf{v}_i))}{\sum_{j=1}^n\exp(s(\mathbf{q},\mathbf{v}_j))}.
\end{equation}

\paragraph{Note on terminology.}
We use ``attention'' as a unifying design label for a family of feature extractors that implement (i) learned reweighting/gating and (ii) structured fusion of sensor features and (optionally) time context. Some extractors are ``attention-inspired'' rather than textbook token-wise attention over sequences; where the implementation deviates from standard formulations, we state this explicitly.

\subsection{Nadaraya--Watson attention (NW)}
NW attention can be interpreted as kernel-regression-based soft retrieval from learnable prototypes. A standard Gaussian-kernel form is
\begin{equation}
\alpha_i = \frac{\exp\left(-\tfrac{\lVert \mathbf{k}_i-\mathbf{q}\rVert^2}{2\sigma^2}\right)}{\sum_j \exp\left(-\tfrac{\lVert \mathbf{k}_j-\mathbf{q}\rVert^2}{2\sigma^2}\right)},
\qquad \mathbf{z}=\sum_i\alpha_i\,\mathbf{v}_i.
\end{equation}
This mechanism is computationally light and aligns with prior PdM attention work motivated by edge constraints \parencite{Siraskar2024NWAttention}.

\paragraph{Theory-to-code mapping.}
In \texttt{code/rl\_pdm.py}, NW attention is implemented as a learned prototype memory with a strictly-positive bandwidth parameter $\beta=\exp(\log\beta)$, which corresponds to $\beta\approx 1/(2\sigma^2)$:
\begin{lstlisting}[language=Python]
# NadarayaWatsonExtractor.forward (code/rl_pdm.py)
query = self.query_net(observations)          # q
dist  = (query.unsqueeze(1) - self.keys.unsqueeze(0)).pow(2).sum(dim=2)
beta  = th.exp(self.log_beta)                 # beta > 0
alpha = th.softmax(-beta * dist, dim=1)       # softmax over prototypes
context = (alpha.unsqueeze(2) * self.values.unsqueeze(0)).sum(dim=1)
return context
\end{lstlisting}

\subsection{Temporal attention (TP)}
Temporal attention injects an explicit stage-of-life signal using sinusoidal positional encodings and a learnable temporal embedding \parencite{Vaswani2017Attention}. It fuses (i) a soft spatial mask over sensor channels and (ii) a time-dependent encoding. This is a lightweight alternative to recurrent memory while still providing a monotonic time context.

\paragraph{Theory-to-code mapping.}
The implementation combines an MLP-based spatial attention mask with a positional encoding lookup indexed by an internal timestep counter:
\begin{lstlisting}[language=Python]
# TemporalAttentionExtractor.forward (code/rl_pdm.py)
spatial_weights = self.spatial_attention(observations)
weighted_spatial = observations * spatial_weights

timestep_idx = int(self.current_timestep.item()) % self.max_timesteps
temporal_encoding = self.positional_encoding[timestep_idx]
temporal_features = temporal_encoding + self.temporal_embedding

combined = th.cat([weighted_spatial, temporal_features.expand(batch_size, -1)], dim=1)
output = self.fusion_net(combined)
self.current_timestep += 1
return output
\end{lstlisting}

\subsection{Multi-head attention (MH)}
Multi-head attention uses multiple attention heads to represent different subspaces of sensor relationships \parencite{Vaswani2017Attention}. In standard form,
\begin{equation}
\mathrm{MultiHead}(Q,K,V)=\mathrm{Concat}(\mathrm{head}_1,\ldots,\mathrm{head}_H)W^O,
\qquad \mathrm{head}_h=\mathrm{Attention}(QW_h^Q,KW_h^K,VW_h^V).
\end{equation}
In our implementation, attention is computed across heads for each observation vector, resulting in a gated mixture of subspace features. Empirically, MH emerges as the strongest and most consistent attention mechanism across datasets (\cref{sec:results}).

\paragraph{Theory-to-code mapping.}
Because the input is a single feature vector (not a token sequence), the code computes a per-head compatibility score $s_h \propto \langle \mathbf{Q}_h,\mathbf{K}_h\rangle$ and normalizes across heads ($h=1,\dots,H$):
\begin{lstlisting}[language=Python]
# MultiHeadAttentionExtractor.forward (code/rl_pdm.py)
Q = self.query_proj(observations).view(batch, H, d)
K = self.key_proj(observations).view(batch, H, d)
V = self.value_proj(observations).view(batch, H, d)

scores  = th.einsum('bhd,bhd->bh', Q, K) * self.scale
alpha_h = th.softmax(scores, dim=1).unsqueeze(2)    # softmax over heads
attended = alpha_h * V
concat = attended.view(batch, H * d)
return self.output_proj(concat)
\end{lstlisting}

\subsection{Self-attention (SA)}
Self-attention projects the observation into a feature space and uses a compact Transformer-style block with residual connections and layer normalization. In standard token-wise form,
\begin{equation}
\mathrm{SelfAttn}(X)=\mathrm{softmax}\left(\frac{XW_Q(XW_K)^\top}{\sqrt{d}}\right)XW_V.
\end{equation}
In this codebase, the softmax normalization is across the \emph{batch} dimension; this can act as a stabilization/gating mechanism, though it differs from token-wise self-attention.

\paragraph{Theory-to-code mapping and deviation from standard self-attention.}
Standard self-attention computes pairwise weights across tokens/features. Here, the implementation computes a scalar score per sample and normalizes across the batch. Concretely, for each batch element $i$ it forms $s_i=(\mathbf{q}_i\odot\mathbf{k}_i)^\top\mathbf{1}/\sqrt{d}$, then $\alpha=\mathrm{softmax}(\mathbf{s})$ across $i$ (batch axis), and outputs $\alpha_i\mathbf{v}_i$ with residual + FFN:
\begin{lstlisting}[language=Python]
# SelfAttentionExtractor.forward (code/rl_pdm.py)
x = self.input_proj(observations)
Q = self.query(x); K = self.key(x); V = self.value(x)

scores = (Q * K).sum(dim=1, keepdim=True) * self.scale
alpha  = th.softmax(scores, dim=0)          # normalize across batch
attn_output = alpha * V

x = self.norm1(x + attn_output)
x = self.norm2(x + self.ffn(x))
return x
\end{lstlisting}

\section{Implementation details and reproducibility}
\label{sec:implementation}
\subsection{REINFORCE implementation}
\label{sec:reinforce_impl}
REINFORCE is implemented as a Stable-Baselines3 compatible \texttt{OnPolicyAlgorithm} (\texttt{code/rl\_pdm.py}). The training loop uses Monte-Carlo returns from the rollout buffer. With an optional value-function baseline, the advantage is $\hat{A}=G-V(s)$, and the policy gradient loss is
\begin{equation}
	\mathcal{L}_{\pi} = -\mathbb{E}[\log\pi(a\mid s)\,\hat{A}].
\end{equation}
This design preserves algorithmic simplicity while still benefiting from variance reduction.

\subsection{Baselines and tooling}
Baselines A2C, DQN, and PPO are accessed via Stable-Baselines3 wrappers and trained in the same environment under the same evaluation protocol \parencite{StableBaselines3}. Attention mechanisms are provided as custom feature extractors.

\section{Experimental design}
\label{sec:experiment}
\subsection{Datasets and generalization protocol}
The study uses two sensor schemas:
\begin{itemize}
	\item \textbf{SIT:} a production-oriented schema with eight sensor channels (vibration, sound, load cells, current).
	\item \textbf{IEEE:} a benchmark schema with force, vibration, and acoustic emission channels.
\end{itemize}

Generalization is evaluated via a \emph{one-vs-all} protocol: agents are trained on one (or a small set of) training file(s) within a schema and evaluated on the remaining unseen files within the same schema.

\subsection{Evaluation metrics}
\label{sec:metrics}
We report three metrics used by the AutoRL evaluation:
\begin{itemize}
	\item \textbf{Evaluation score:} a weighted score combining tool usage, $\lambda$, and threshold violations.
	\item \textbf{Tool usage (\%):} percentage of wear threshold utilized at first replacement; higher is better.
	\item \textbf{$\lambda$ timing metric:} a signed timing error between desired first replacement timing and the actual first replacement time.
\end{itemize}

\paragraph{Exact evaluation score definition (as implemented).}
The paper-level ``evaluation score'' corresponds to \texttt{calculate\_eval\_score()} in \texttt{code/batch\_train.py}. Let $u\in[0,1]$ be tool usage fraction at first replacement (or a penalized proxy if no replacement occurs), $v\in\mathbb{N}_0$ be the number of post-threshold timesteps without replacement (threshold violations), and $\lambda=t_{\lambda}-t_{\mathrm{FR}}$ be the signed timing error (\cref{sec:metrics}). The implemented score is
\begin{align}
S_{\text{tool}} &= \mathrm{clip}(u,0,1),\\
S_{\lambda} &= \max\left(0,\,1-\frac{|\lambda|}{c_{\lambda}}\right),\\
S_{\text{viol}} &= \frac{1}{1+v},\\
\mathrm{EvalScore} &= 0.50\,S_{\text{tool}} + 0.30\,S_{\lambda} + 0.20\,S_{\text{viol}}.
\end{align}
Here $c_{\lambda}$ is the normalization constant used by \texttt{calculate\_eval\_score()}, implemented as $c_{\lambda}=\texttt{WEAR\_THRESHOLD}\,(1-\texttt{IAR\_RANGE})$. Note that in the current codebase $\lambda$ is measured in \emph{timesteps} (\texttt{rl\_pdm.evaluate\_trained\_model}), whereas $c_{\lambda}$ is derived from a \emph{wear} threshold; thus the ratio $|\lambda|/c_{\lambda}$ should be interpreted as an implementation-defined heuristic scaling rather than a dimensionally normalized quantity. Additionally, the $t_{\lambda}$ event is triggered using \texttt{IAR\_lower} computed in \texttt{rl\_pdm.py} as $\texttt{WEAR\_THRESHOLD}-2\,(\texttt{WEAR\_THRESHOLD}\cdot\texttt{IAR\_RANGE})$, which differs from $c_{\lambda}$.

\paragraph{Implementation reference.}
\begin{lstlisting}[language=Python]
# calculate_eval_score (code/batch_train.py)
tool_usage_score = min(1.0, max(0.0, tool_usage_pct))
IAR_LOWER = rl_pdm.WEAR_THRESHOLD * (1.0 - rl_pdm.IAR_RANGE)
lambda_score = max(0.0, 1.0 - (abs(lambda_metric) / IAR_LOWER))
violation_score = max(0.0, 1.0 / (1.0 + violations))
eval_score = 0.50*tool_usage_score + 0.30*lambda_score + 0.20*violation_score
\end{lstlisting}

\paragraph{$\lambda$ metric and offline prognostic evaluation.}
We implement a timing metric inspired by offline prognostic evaluation concepts \parencite{Saxena2008Metrics}. Define $t_{\lambda}$ as the first timestep where wear exceeds an ``in-advance replacement'' (IAR) lower bound. Let $t_{\mathrm{FR}}$ be the first replacement timestep. Then
\begin{equation}
	\lambda = t_{\lambda} - t_{\mathrm{FR}}.
\end{equation}
Lower $|\lambda|$ is better; $\lambda=0$ corresponds to replacing exactly at the desired timing point. In the implementation, IAR bounds are derived from the wear threshold using a configurable range parameter.

\subsection{Multi-round evaluation, confidence intervals, and hypothesis tests}
Each trained model is evaluated for 20 rounds on unseen data. We report mean values and 95\% confidence intervals where raw results are available (IEEE). Hypothesis testing uses one-tailed Welch $t$-tests (provided as CSV artefacts) at $\alpha=0.05$.

\paragraph{Aggregation and 95\% confidence intervals.}
Let $S_1,\ldots,S_n$ denote the evaluation scores obtained across rounds and test files in a given panel. We report the sample mean and standard deviation
\begin{align}
\mu &= \frac{1}{n}\sum_{i=1}^{n} S_i,\\
\sigma &= \sqrt{\frac{1}{n-1}\sum_{i=1}^{n}(S_i-\mu)^2},
\end{align}
and the standard error of the mean $\mathrm{SEM}=\sigma/\sqrt{n}$. A normal-approximation 95\% CI is then $[\mu-1.96\,\mathrm{SEM},\;\mu+1.96\,\mathrm{SEM}]$.

\paragraph{What varies across rounds (policy stochasticity + seeding).}
Multi-round evaluation is designed to quantify variance from \emph{stochastic decision policies} rather than resampling the dataset (the environment replays a fixed run). Concretely, \texttt{code/batch\_train.py} assigns a distinct seed per (round, retry), and calls \texttt{evaluate\_trained\_model(..., deterministic=False)} so actions are sampled stochastically from the policy. To keep temporal feature extractors consistent with episode progression, \texttt{evaluate\_trained\_model} executes predictions step-by-step (not as a single batch), and attempts to reset any internal timestep counters at the start of evaluation.

\paragraph{Statistical caveats.}
The $t$-tests are used as a practical, automated comparison tool for the AutoRL grid and should be interpreted with standard cautions: (i) multiple comparisons are performed across algorithms/attention panels and no correction is applied, and (ii) repeated evaluation rounds on the same test file can introduce dependence (pseudoreplication). Accordingly, the statistical panels are best treated as evidence of consistent trends rather than as definitive confirmatory inference; effect sizes and robustness checks (e.g., bootstrap confidence intervals or nonparametric tests) are recommended in follow-up analysis.

\subsection{Training budget and default configuration}
The batch pipeline enumerates learning rates and discount factors ($\gamma$) as part of the AutoRL grid (stored in the checkpoint queue and reported with the analysis artefacts). For reproducibility, key environment constants (thresholds, reward weights, IAR range, and default learning settings) are summarized in \cref{app:defaults}. These constants are taken directly from the implementation in \texttt{code/rl\_pdm.py}.

\subsection{Panels and reporting}
The hypothesis-test artefacts provide multiple panels (e.g., ``with self-eval'' vs ``unseen only'', and ``all attention'' vs ``best attention''). This paper emphasizes \emph{unseen-only} panels because they align with the intended PdM deployment assumption: policies should generalize to new runs rather than overfit to the training trace.

% -----------------------------------------------------------------------------
% Results (kept essentially intact from AutoRL_for_PdM.tex)
% -----------------------------------------------------------------------------
\section{Results}
\label{sec:results}
This section presents results for SIT Set-1, SIT Set-2, and IEEE using the provided analysis dashboards, heatmaps, and statistical heatmaps. IEEE additionally includes an evaluation CSV enabling summary tables.

\subsection{Visual summaries: analysis dashboards and generalization heatmaps}
\Cref{fig:analysis_reports} presents consolidated analysis dashboards for each schema/set, including (i) overall performance matrices, (ii) CI-based error bars, (iii) algorithm-level aggregates, and (iv) $\lambda$ behavior. \Cref{fig:heatmaps} shows one-vs-all generalization heatmaps.

\begin{figure}[t]
	\centering
	\begin{subfigure}[t]{0.32\linewidth}
		\centering
		\includegraphics[width=\linewidth]{\autorlFigRoot/SIT_Results/_Analysis_Report_SIT_Set-1.jpg}
		\caption{SIT Set-1}
		\label{fig:analysis_sit1}
	\end{subfigure}
	\begin{subfigure}[t]{0.32\linewidth}
		\centering
		\includegraphics[width=\linewidth]{\autorlFigRoot/SIT_Results/_Analysis_Report_SIT_Set-2.jpg}
		\caption{SIT Set-2}
		\label{fig:analysis_sit2}
	\end{subfigure}
	\begin{subfigure}[t]{0.32\linewidth}
		\centering
		\includegraphics[width=\linewidth]{\autorlFigRoot/IEEE_Results/_Analysis_Report_IEEE.jpg}
		\caption{IEEE}
		\label{fig:analysis_ieee}
	\end{subfigure}
	\caption{Consolidated evaluation dashboards (mean and CI error bars), used as primary visual evidence for performance patterns across algorithms and attention mechanisms.}
	\label{fig:analysis_reports}
\end{figure}

\begin{figure}[t]
	\centering
	\begin{subfigure}[t]{0.32\linewidth}
		\centering
		\includegraphics[width=\linewidth]{\autorlFigRoot/SIT_Results/_Heatmap_SIT_Set-1.jpg}
		\caption{SIT Set-1}
	\end{subfigure}
	\begin{subfigure}[t]{0.32\linewidth}
		\centering
		\includegraphics[width=\linewidth]{\autorlFigRoot/SIT_Results/_Heatmap_SIT_Set-2.jpg}
		\caption{SIT Set-2}
	\end{subfigure}
	\begin{subfigure}[t]{0.32\linewidth}
		\centering
		\includegraphics[width=\linewidth]{\autorlFigRoot/IEEE_Results/_Heatmap_IEEE.jpg}
		\caption{IEEE}
	\end{subfigure}
	\caption{One-vs-all generalization heatmaps: trained models (rows) evaluated on unseen test files (columns) within each schema.}
	\label{fig:heatmaps}
\end{figure}

\subsection{Hypothesis tests (H1): REINFORCE vs A2C/DQN/PPO}
\label{sec:hypothesis_tests}
The hypothesis-test CSV artefacts provide one-tailed Welch $t$-tests at $\alpha=0.05$. \Cref{tab:hyp_tests} summarizes the unseen-only evaluation-score tests for two panels: \emph{All Attn} (aggregating REINFORCE across attention mechanisms) and \emph{Best Attn} (restricting REINFORCE to the best attention mechanism for that dataset panel).

\begin{table}[t]
	\centering
	\caption{Hypothesis tests for H1 (unseen-only evaluation score): one-tailed Welch $t$-tests comparing REINFORCE to each competitor at $\alpha=0.05$.}
	\label{tab:hyp_tests}
	% Auto-generated from *_Hypothesis_Tests_*.csv (Eval Score, Unseen Only)
\begin{tabular}{l l l r r l}
\toprule
\rowcolor{gray!20} \textbf{Dataset} & \textbf{Panel} & \textbf{Competitor} & \textbf{t} & \textbf{p (1-tail)} & \textbf{Sig.}\\
\midrule
SIT & All Attn & A2C & 198.007 & 0.000 & YES\\
SIT & Best Attn & A2C & 512.423 & 0.000 & YES\\
SIT & All Attn & DQN & 202.449 & 0.000 & YES\\
SIT & Best Attn & DQN & 562.313 & 0.000 & YES\\
SIT & All Attn & PPO & 20.686 & 0.000 & YES\\
SIT & Best Attn & PPO & 49.902 & 0.000 & YES\\
IEEE & All Attn & A2C & 10.726 & 0.000 & YES\\
IEEE & Best Attn & A2C & 10.434 & 0.000 & YES\\
IEEE & All Attn & DQN & 15.977 & 0.000 & YES\\
IEEE & Best Attn & DQN & 15.235 & 0.000 & YES\\
IEEE & All Attn & PPO & 1.957 & 0.026 & YES\\
IEEE & Best Attn & PPO & 2.096 & 0.020 & YES\\
\bottomrule
\end{tabular}

\end{table}

Across SIT Set-1, SIT Set-2, and IEEE, the unseen-only tests show strong evidence that REINFORCE outperforms A2C and DQN. Against PPO, the evidence is schema-dependent: it is significant for SIT Set-1 and IEEE, and not significant in SIT Set-2 under the aggregated unseen-only panel. This pattern motivates H3: attention mechanisms can help differentiate configurations where baseline algorithms are otherwise close.

\subsection{Dataset-level interpretation}
\label{sec:dataset_interpretation}
\paragraph{SIT Set-1.}
SIT Set-1 dashboards and heatmaps indicate strong performance for REINFORCE across attention mechanisms, with multi-head attention improving consistency. The one-vs-all heatmap also suggests that some training traces transfer better than others, motivating a multi-round, multi-test-file evaluation rather than a single split.

\paragraph{SIT Set-2.}
SIT Set-2 is more challenging for the PPO comparison. In the unseen-only \emph{All Attn} panel, REINFORCE does not significantly outperform PPO, but the \emph{Best Attn} panel does show significance. Practically, this means representation choice (attention) is sometimes the decisive factor and supports H3 as an applied claim: a lightweight algorithm with a suitable feature extractor can match or exceed a heavier baseline.

\paragraph{IEEE.}
IEEE results show that top configurations cluster around PPO-Temporal and REINFORCE variants with multi-head/self-attention/temporal attention (\cref{tab:ieee_top10}). This clustering suggests that (i) partial observability and sensor fusion dominate difficulty, and (ii) attention-driven representation learning is central to performance.

\subsection{Quantitative IEEE summary (means and 95\% CI)}
\label{sec:ieee_summary}
For IEEE, a raw evaluation CSV enables direct computation of means and 95\% confidence intervals on unseen-only evaluations. \Cref{tab:ieee_top10} lists the top configurations by mean evaluation score, and \Cref{tab:ieee_best_attn} summarizes the best attention mechanism per algorithm.

\begin{table}[t]
	\centering
	\caption{Top IEEE configurations on unseen-only evaluations (mean and 95\% CI, aggregated across rounds and test files). Numbers are rounded to three decimals.}
	\label{tab:ieee_top10}
	\input{\autorlTableRoot/ieee_unseen_summary.tex}
\end{table}

\begin{table}[t]
	\centering
	\caption{IEEE unseen-only: best attention mechanism per algorithm (mean evaluation score and 95\% CI).}
	\label{tab:ieee_best_attn}
	\input{\autorlTableRoot/ieee_best_attention_per_algo.tex}
\end{table}

\subsection{Attention improves REINFORCE (H3)}
Across schemas, multi-head attention emerges as the best or most consistent attention mechanism in the statistical panels and the consolidated analysis reports (\cref{fig:analysis_reports}). From a PdM perspective, this is plausible because multi-head gating can capture heterogeneous sensor relationships (e.g., load-driven vs vibration-driven wear regimes) without requiring a large recurrent model.

\subsection{Practical recommendation}
For edge-first deployments where training and inference budgets are constrained, the evidence supports a simple default: \emph{REINFORCE + Multi-Head attention}. When a practitioner has bandwidth to benchmark alternatives, PPO with temporal attention is a strong secondary candidate (particularly on IEEE).

\subsection{Statistical heatmaps: p-values and t-statistics}
\begin{figure}[t]
	\centering
	\begin{subfigure}[t]{0.32\linewidth}
		\centering
		\includegraphics[width=\linewidth]{\autorlFigRoot/SIT_Results/_Statistical_p-value_SIT_Set-1.jpg}
		\caption{SIT Set-1}
	\end{subfigure}
	\begin{subfigure}[t]{0.32\linewidth}
		\centering
		\includegraphics[width=\linewidth]{\autorlFigRoot/SIT_Results/_Statistical_p-value_SIT_Set-2.jpg}
		\caption{SIT Set-2}
	\end{subfigure}
	\begin{subfigure}[t]{0.32\linewidth}
		\centering
		\includegraphics[width=\linewidth]{\autorlFigRoot/IEEE_Results/_Statistical_p-value_IEEE.jpg}
		\caption{IEEE}
	\end{subfigure}
	\caption{$p$-value heatmaps for hypothesis tests (blue indicates stronger evidence).}
	\label{fig:pvalues}
\end{figure}

\begin{figure}[t]
	\centering
	\begin{subfigure}[t]{0.32\linewidth}
		\centering
		\includegraphics[width=\linewidth]{\autorlFigRoot/SIT_Results/_Statistical_t_SIT_Set-1.jpg}
		\caption{SIT Set-1}
	\end{subfigure}
	\begin{subfigure}[t]{0.32\linewidth}
		\centering
		\includegraphics[width=\linewidth]{\autorlFigRoot/SIT_Results/_Statistical_t_SIT_Set-2.jpg}
		\caption{SIT Set-2}
	\end{subfigure}
	\begin{subfigure}[t]{0.32\linewidth}
		\centering
		\includegraphics[width=\linewidth]{\autorlFigRoot/IEEE_Results/_Statistical_t_IEEE.jpg}
		\caption{IEEE}
	\end{subfigure}
	\caption{$t$-statistic heatmaps corresponding to \cref{fig:pvalues}.}
	\label{fig:tstats}
\end{figure}

\subsection{H2: cross-schema generality (SIT vs IEEE)}
H2 is supported structurally by the environment design (\cref{sec:env}) and empirically by results on both schemas (\cref{fig:analysis_reports,fig:heatmaps}). The same environment class (with schema detection and feature selection) and the same reward shaping (\cref{sec:reward}) produce strong-performing agents in both SIT and IEEE settings.

\section{Discussion}
\label{sec:discussion}
\subsection{Why lightweight REINFORCE can win in this PdM setup}
The PdM task studied here has a small action space and terminates on replacement; accordingly, the control problem can resemble learning a decision boundary over sensor-derived features. When partial observability dominates the difficulty (\cref{sec:pomdp}), representation quality can matter more than optimizer sophistication. Attention mechanisms (especially MH and TP) inject inductive biases that reduce the burden on the RL algorithm itself.

\subsection{Implications for industrial practitioners}
AutoRL reduces reliance on RL expertise by automating algorithm/attention sweeps and exposing standardized result artefacts (dashboards, heatmaps, and hypothesis-test summaries). This directly supports RO3: practitioners can select a robust configuration empirically, without deeply understanding policy-gradient variance, PPO clipping, or DQN stability tricks.

% From UM_bacground:
\subsection{Edge compute and sustainability framing}
This paper does not directly measure energy consumption. However, the combination of AutoRL-style sweeps and edge deployment makes compute-efficiency a meaningful engineering constraint. In this context, if REINFORCE achieves equal or better PdM outcomes than heavier baselines, it supports a defensible ``Green AI'' argument: prefer the simplest method that meets the operational objective \parencite{Schwartz2020GreenAI, Strubell2019Energy}.

\section{Limitations and threats to validity}
\label{sec:limits}
\begin{itemize}
	\item \textbf{Dataset-driven environment:} the environment replays logged runs-to-threshold and does not model intervention dynamics beyond termination; future work can extend to maintenance logistics and downtime costs.
	\item \textbf{Temporal extractor state:} temporal attention maintains an internal timestep counter; production use should ensure the counter resets consistently with episode resets.
	\item \textbf{Compute and carbon measurement:} energy/carbon implications are argued from algorithmic complexity and deployment constraints rather than measured on edge hardware.
\end{itemize}

\section{Conclusions and future work}
\label{sec:conclusion}
This paper presented an AutoRL pipeline and a schema-adaptive Gymnasium PdM environment for milling-machine tool replacement. Empirical results across SIT and IEEE schemas, supported by hypothesis tests at $\alpha=0.05$, show that REINFORCE is a strong candidate for edge-first PdM, and that attention mechanisms---especially multi-head attention---can further improve performance and robustness.

Future work should (i) measure inference latency and energy on representative edge devices, (ii) incorporate richer maintenance costs and partial observability models, and (iii) evaluate robustness under sensor drift and cross-machine domain shift.

\appendix
\section{Environment and reward details}
\label{app:env_reward}
This appendix summarizes implementation-facing details to support reproducibility and practitioner adoption.

\subsection{Observation design and POMDP rationale}
The environment returns observation vectors comprised only of sensor channels. Tool wear is excluded from observations but is used internally to compute termination and reward. This enforces partial observability and prevents leakage of the health label into the policy input.

\subsection{Reward components}
The reward logic uses constants $R_1$--$R_4$ and a large catastrophic penalty when the wear threshold is violated. Replacement receives a cost and, when performed near the threshold, an additional bonus. \Cref{tab:reward_components} presents a qualitative summary.

\begin{table}[t]
	\centering
	\caption{Reward shaping components in the milling PdM environment (qualitative summary).}
	\label{tab:reward_components}
	\begin{tabular}{p{0.20\linewidth} p{0.22\linewidth} p{0.50\linewidth}}
		\toprule
		\rowcolor{gray!20} \textbf{Component} & \textbf{Applied when} & \textbf{Effect}\\
		\midrule
		$R_1$ survival reward & continue and safe & Encourages utilization while below threshold.\\
		$R_2$ violation penalty & continue and violate & Strongly discourages threshold violation.\\
		$R_3$ replacement cost & replace action & Regularizes against replacing too early.\\
		$R_4$ near-threshold bonus & replace near $\tau$ & Rewards timely replacement close to threshold.\\
		Stepwise bonuses & wear ratio $\,\ge 0.7, 0.8, 0.9$ & Increases incentive as wear approaches threshold.\\
		\bottomrule
	\end{tabular}
\end{table}

\section{Evaluation protocol details}
\label{app:eval}
\subsection{Multi-round evaluation}
Each agent is evaluated for multiple rounds per test file to reduce variance due to stochastic policies and initialization. For IEEE, this enables confidence intervals on unseen-only evaluations.

\subsection{Unseen-only emphasis}
This paper focuses on unseen-only evaluation panels for generalization claims (H2). Panels that include self-evaluation are useful for debugging and sanity checks but can overestimate deployment performance.

\section{Dataset schemas and sensor signals}
\label{app:schema}
The environment supports two primary schemas through column-based feature selection. \Cref{tab:schema_features} lists the sensor channels used in each schema (excluding \texttt{Time}, \texttt{tool\_wear}, and \texttt{ACTION\_CODE}).

\begin{table}[t]
\centering
\caption{Sensor channels used as observations by schema (as implemented in \texttt{MT\_Env}).}
\label{tab:schema_features}
\begin{tabular}{p{0.18\linewidth} p{0.75\linewidth}}
\toprule
\rowcolor{gray!20} \textbf{Schema} & \textbf{Observation channels}\\
\midrule
SIT & \texttt{Vib\_Spindle}, \texttt{Vib\_Table}, \texttt{Sound\_Spindle}, \texttt{Sound\_table}, \texttt{X\_Load\_Cell}, \texttt{Y\_Load\_Cell}, \texttt{Z\_Load\_Cell}, \texttt{Current}\\
IEEE & \texttt{force\_x}, \texttt{force\_y}, \texttt{force\_z}, \texttt{vibration\_x}, \texttt{vibration\_y}, \texttt{vibration\_z}, \texttt{acoustic\_emission\_rms}\\
\bottomrule
\end{tabular}
\end{table}

\section{Default constants and hyperparameters}
\label{app:defaults}
\Cref{tab:defaults} summarizes the default constants used in the PdM environment and training module. AutoRL sweeps may override some values (e.g., learning rate, $\gamma$) via the CLI pipeline.

\begin{table}[t]
\centering
\caption{Default environment/training constants (from \texttt{code/rl\_pdm.py} and \texttt{code/batch\_train.py}).}
\label{tab:defaults}
\begin{tabular}{p{0.34\linewidth} p{0.18\linewidth} p{0.40\linewidth}}
\toprule
\rowcolor{gray!20} \textbf{Parameter} & \textbf{Default} & \textbf{Meaning}\\
\midrule
\texttt{WEAR\_THRESHOLD} & 300.000 & Wear threshold used for termination and reporting.\\
\texttt{TRAINING\_WEAR\_THRESHOLD} & 300.000 & Training threshold (configurable; default equals threshold).\\
\texttt{EPISODES} & 100 & Training episodes per run (default in RL module).\\
\texttt{LR\_DEFAULT} & 0.001 & Default learning rate (overridden by AutoRL sweep values).\\
\texttt{GAMMA\_DEFAULT} & 0.990 & Default discount factor (overridden by AutoRL sweep values).\\
\texttt{EVAL\_ROUNDS} & 20 & Evaluation rounds per model on unseen data.\\
\texttt{IAR\_RANGE} & 0.050 & IAR bounds range for timing metric computation.\\
\texttt{LAMBDA} & 0.900 & Logistic slack coefficient used in $\lambda$ timing computation.\\
\texttt{R1}, \texttt{R2}, \texttt{R3}, \texttt{R4} & 0.100, -100.000, -0.500, 60.000 & Reward shaping constants.\\
\bottomrule
\end{tabular}
\end{table}

\printbibliography[title={References}]

\end{document}
